\documentclass[12pt,a4paper]{article}
% !TEX root = ../report_plan/report_draft/RENJIEQU_235_FinalReport.tex
\usepackage[utf8]{inputenc}
\usepackage{amsmath,amssymb,amsthm}
\usepackage{geometry}
\usepackage{hyperref}
\usepackage{enumitem}
\usepackage{titlesec}
\usepackage{booktabs}

\geometry{margin=2.5cm}
\hypersetup{colorlinks=true, linkcolor=blue, citecolor=blue, urlcolor=blue}

% --- robust bra-ket and operator macros (avoid package version clashes) ---
\newcommand{\ket}[1]{\lvert #1 \rangle}
\newcommand{\bra}[1]{\langle #1 \rvert}
\newcommand{\braket}[2]{\langle #1 \,|\, #2 \rangle}
\newcommand{\mel}[3]{\langle #1 \,|\, #2 \,|\, #3 \rangle}
\newcommand{\re}{\operatorname{Re}}
\newcommand{\im}{\operatorname{Im}}
\newcommand{\dlam}{\dot{\boldsymbol{\lambda}}}
\newcommand{\Hexp}{\langle H\rangle}

\newtheorem{definition}{Definition}[section]
\newtheorem{theorem}{Theorem}[section]
\newtheorem{proposition}{Proposition}[section]
\newtheorem{remark}{Remark}[section]
\newtheorem{example}{Example}[section]

\title{\textbf{Variational Quantum Simulation of Real-Time Dynamics}\\[4pt]
\large McLachlan's Variational Principle: Equations of Motion, the Global Phase, and the Geometric Projector}
\author{Theory notes}
\date{\today}

\begin{document}
\maketitle
\tableofcontents
\bigskip

\noindent\textbf{Scope.} These notes treat \emph{only} the variational quantum simulation (VQS)
algorithm for real-time evolution, built on McLachlan's variational principle (MVP). We derive
the equations of motion from first principles, explain the role of the energy-variance term and of
the global phase, show how the geometric projector $Q_\psi$ emerges, lay out the iteration loop,
and analyse two minimal models that expose when the algorithm succeeds and fails. The application
context is the quench dynamics of the lattice Schwinger model (initial state from VQE, evolution
under a quenched Hamiltonian), but everything below is independent of that specific Hamiltonian.

\newpage
\part{Variational Principle and Equations of Motion}

\section{Setup: the real-time problem on a variational manifold}\label{sec:setup}

We wish to approximate the solution of the time-dependent Schr\"odinger equation (units $\hbar=1$),
\begin{equation}
  \frac{d}{dt}\ket{\Psi(t)} = -iH\ket{\Psi(t)},
  \label{eq:schrodinger}
\end{equation}
without compiling the propagator $e^{-iHt}$ into a circuit whose depth grows with $t$ (the
Suzuki--Trotter strategy). Instead we restrict the state to a \emph{fixed-depth variational
manifold}
\begin{equation}
  \ket{\psi(\boldsymbol{\lambda})}=U(\boldsymbol{\lambda})\ket{\psi_0},
  \qquad \boldsymbol{\lambda}=(\lambda_1,\dots,\lambda_{N_p})\in\mathbb{R}^{N_p},
  \label{eq:ansatz}
\end{equation}
and encode the time dependence \emph{entirely} through the parameters, $\boldsymbol{\lambda}=\boldsymbol{\lambda}(t)$.
The problem of solving \eqref{eq:schrodinger} then reduces to finding a trajectory
$\boldsymbol{\lambda}(t)$ such that $\ket{\psi(\boldsymbol{\lambda}(t))}$ best tracks the exact state.

\begin{definition}[Tangent vectors]
At a point $\boldsymbol{\lambda}$ on the manifold, the \emph{parameter derivatives}
\[
  \ket{\partial_i\psi}\equiv\frac{\partial}{\partial\lambda_i}\ket{\psi(\boldsymbol{\lambda})}
\]
span the tangent space. Because the state depends on time only through $\boldsymbol{\lambda}(t)$,
the chain rule gives the manifold velocity
\begin{equation}
  \ket{\dot\psi}=\sum_{i=1}^{N_p}\dot\lambda_i\ket{\partial_i\psi}.
  \label{eq:vel}
\end{equation}
The unknowns are the \emph{real} velocities $\dot\lambda_i$; the vectors $\ket{\partial_i\psi}$ and
$\ket{\psi}$ are fixed data at the current $\boldsymbol{\lambda}$.
\end{definition}

Throughout we assume normalization $\braket{\psi}{\psi}=1$, from which
\begin{equation}
  \partial_i\braket{\psi}{\psi}=\braket{\partial_i\psi}{\psi}+\braket{\psi}{\partial_i\psi}=0
  \;\Longrightarrow\;
  \re\braket{\psi}{\partial_i\psi}=0,
  \label{eq:normid}
\end{equation}
i.e.\ $\braket{\psi}{\partial_i\psi}$ is purely imaginary. We will use this repeatedly.

\section{McLachlan's variational principle}\label{sec:mvp}

\begin{definition}[McLachlan distance]
The \emph{McLachlan functional} is the squared $\ell^2$ residual between the variational
velocity and the exact velocity $-iH\ket{\psi}$,
\begin{equation}
  L(\dlam)\equiv\Big\|\,\ket{\dot\psi}+iH\ket{\psi}\,\Big\|^2 .
  \label{eq:Lfunc}
\end{equation}
\end{definition}

The principle is: \emph{at each instant, choose the tangent vector $\ket{\dot\psi}$ (equivalently the
velocities $\dlam$) that minimizes $L$}. Geometrically this is the orthogonal projection of the exact
velocity onto the tangent space. The variation is taken with respect to $\dlam$ at \emph{fixed}
$\boldsymbol{\lambda}$ (hence at fixed $\ket{\psi}$ and $\ket{\partial_i\psi}$).

\section{Expansion of the functional}\label{sec:expand}

Using $\|a+b\|^2=\|a\|^2+2\re\braket{a}{b}+\|b\|^2$ with $a=\ket{\dot\psi}$ and $b=iH\ket{\psi}$,
\begin{equation}
  L=\braket{\dot\psi}{\dot\psi}+2\re\!\big(i\mel{\dot\psi}{H}{\psi}\big)+\mel{\psi}{H^2}{\psi}.
  \label{eq:Lexp1}
\end{equation}
The cross term is where an imaginary part appears: since $\re(iz)=-\im(z)$ for any complex $z$,
\begin{equation}
  2\re\!\big(i\mel{\dot\psi}{H}{\psi}\big)=-2\im\mel{\dot\psi}{H}{\psi},
  \label{eq:crossterm}
\end{equation}
so that
\begin{equation}
  L=\underbrace{\braket{\dot\psi}{\dot\psi}}_{\text{quadratic in }\dlam}
   \;-\;\underbrace{2\,\im\mel{\dot\psi}{H}{\psi}}_{\text{linear in }\dlam}
   \;+\;\underbrace{\langle H^2\rangle}_{\text{constant}},
  \label{eq:Lexp2}
\end{equation}
where $\langle H^2\rangle\equiv\mel{\psi}{H^2}{\psi}$.

\begin{remark}[Why $\langle H^2\rangle$ is ``constant'']
The state $\ket{\psi}$ certainly changes along the trajectory, so $\langle H^2\rangle$ changes
\emph{in time}. But within a single minimization step it does not depend on the optimization
variable $\dlam$: the term contains $\ket{\psi}$ only, never $\ket{\dot\psi}$. Hence
$\partial\langle H^2\rangle/\partial\dot\lambda_k=0$ and it drops out of the stationarity condition.
This is the same reason the additive constant $c$ in $f(v)=\tfrac12 a v^2-bv+c$ never affects the
minimizer $v^\ast=b/a$: it shifts the parabola vertically without moving its vertex.
\end{remark}

Substituting the tangent expansion \eqref{eq:vel} into \eqref{eq:Lexp2}, with $\dot\lambda_i$ real:
\begin{align}
  \braket{\dot\psi}{\dot\psi}
    &=\sum_{i,j}\dot\lambda_i\dot\lambda_j\,\braket{\partial_i\psi}{\partial_j\psi}
     =\sum_{i,j}\dot\lambda_i\dot\lambda_j\,\re\braket{\partial_i\psi}{\partial_j\psi},
     \label{eq:quad}\\[2pt]
  -2\,\im\mel{\dot\psi}{H}{\psi}
    &=-2\sum_i\dot\lambda_i\,\im\mel{\partial_i\psi}{H}{\psi}.
     \label{eq:lin}
\end{align}
In \eqref{eq:quad} only the real part survives because
$\im\braket{\partial_i\psi}{\partial_j\psi}$ is antisymmetric in $i\leftrightarrow j$
(as $\braket{\partial_j\psi}{\partial_i\psi}=\overline{\braket{\partial_i\psi}{\partial_j\psi}}$),
while $\dot\lambda_i\dot\lambda_j$ is symmetric, so the antisymmetric piece sums to zero.

\section{The linear system}\label{sec:linsys}

Define the symmetric matrix $M$ and the vector $V$,
\begin{equation}
  M_{ij}\equiv\re\braket{\partial_i\psi}{\partial_j\psi},
  \qquad
  V_i\equiv\im\mel{\partial_i\psi}{H}{\psi}.
  \label{eq:MVdef}
\end{equation}
Then \eqref{eq:Lexp2} is exactly a quadratic form in the velocity,
\begin{equation}
  L(\dlam)=\dlam^{\!\top}M\,\dlam-2\,V^{\!\top}\dlam+\langle H^2\rangle.
  \label{eq:quadform}
\end{equation}
Minimizing a (positive) quadratic form is a linear problem: setting
$\partial L/\partial\dot\lambda_k=0$ gives the \emph{normal equations}
\begin{equation}
  \frac{\partial L}{\partial\dot\lambda_k}
   =2\sum_j M_{kj}\dot\lambda_j-2V_k=0
   \;\Longrightarrow\;
   \boxed{\;\sum_j M_{kj}\,\dot\lambda_j=V_k\;}
  \label{eq:EOM}
\end{equation}
or in matrix form $M\dlam=V$, solved at every instant for the velocity $\dlam=M^{-1}V$.

\begin{remark}[Physical reading of $M$ and $V$]
$M$ is the real part of the quantum geometric tensor (the Fubini--Study metric); it converts
parameter velocities into state velocities. $V$ is a projected energy gradient pointing along the
direction the true Schr\"odinger flow demands. The solution $\dlam=M^{-1}V$ is the best
least-squares image, inside the achievable tangent space, of the exact velocity $-iH\ket{\psi}$.
\end{remark}

\section{The residual and the role of $\langle H^2\rangle$}\label{sec:residual}

Although $\langle H^2\rangle$ does not affect $\dlam$, it does fix the \emph{value} of the residual.
Substituting the optimum $\dlam^\ast=M^{-1}V$ into \eqref{eq:quadform}, and using $M^{\!\top}=M$,
\begin{align}
  L_{\min}
  &=(M^{-1}V)^{\!\top}M(M^{-1}V)-2V^{\!\top}(M^{-1}V)+\langle H^2\rangle\nonumber\\
  &=V^{\!\top}M^{-1}V-2V^{\!\top}M^{-1}V+\langle H^2\rangle\nonumber\\
  &=\boxed{\;\langle H^2\rangle-V^{\!\top}M^{-1}V\;}.
  \label{eq:residual}
\end{align}
$L_{\min}$ is the squared McLachlan distance: how far the best variational velocity remains from
the exact one.

\begin{remark}[Practical consequence]
To determine \emph{which way to step} (the dynamics) one needs only $M$ and $V$; $\langle H^2\rangle$
is unnecessary. To monitor \emph{how well} the ansatz tracks, one needs $L_{\min}$ and hence
$\langle H^2\rangle$. On a small classical emulator $\langle H^2\rangle$ is cheap and $L_{\min}(t)$
is a useful quality diagnostic, complementary to the fidelity against an exact reference.
\end{remark}

\section{Consistency check: exactly representable dynamics}\label{sec:check}

\begin{proposition}
If the exact velocity lies in the tangent space, i.e.\ there exist real $\dot\lambda_i$ with
$\sum_i\dot\lambda_i\ket{\partial_i\psi}=-iH\ket{\psi}$, then \eqref{eq:EOM} reproduces it.
\end{proposition}
\begin{proof}
Insert $\ket{\dot\psi}=-iH\ket{\psi}$ into the left side of \eqref{eq:EOM}:
\[
  \sum_j M_{ij}\dot\lambda_j=\re\!\Big\langle\partial_i\psi\,\Big|\,\sum_j\dot\lambda_j\,\partial_j\psi\Big\rangle
  =\re\mel{\partial_i\psi}{(-iH)}{\psi}.
\]
Using $\re(-iz)=\im(z)$,
\[
  \re\big(-i\mel{\partial_i\psi}{H}{\psi}\big)=\im\mel{\partial_i\psi}{H}{\psi}=V_i,
\]
which matches the right side of \eqref{eq:EOM}. Hence the exact $\dlam$ solves the system. This
also fixes the sign convention in \eqref{eq:MVdef}: $V_i=\im\mel{\partial_i\psi}{H}{\psi}$ (not its
negative); the opposite sign would run time backwards.
\end{proof}

\newpage
\part{The Global Phase and the Geometric Projector}

\section{The global-phase parameter}\label{sec:phase}

A global phase is physically unobservable: $e^{i\lambda_0}\ket{\psi}$ and $\ket{\psi}$ give identical
expectation values. Nonetheless the exact velocity $-iH\ket{\psi}$ generically has a component
\emph{along} $\ket{\psi}$ (a phase rotation) of magnitude $\langle H\rangle$, because
$\braket{\psi}{-iH\psi}=-i\langle H\rangle\neq0$. We therefore augment the ansatz with an explicit
phase knob,
\begin{equation}
  \ket{\psi(\boldsymbol{\lambda})}=e^{i\lambda_0(t)}\prod_{i=1}^{N_p}U(\lambda_i(t))\ket{0},
  \qquad
  \ket{\partial_0\psi}=i\ket{\psi}.
  \label{eq:phaseansatz}
\end{equation}
The extra tangent vector $\ket{\partial_0\psi}=i\ket{\psi}$ points purely along $\ket{\psi}$: it is a
pure phase direction with zero physical content.

\begin{remark}[Why add a knob that does nothing to the state]
The physical state lives in projective Hilbert space (rays). Without a dedicated phase direction,
the minimization \eqref{eq:Lfunc} forces the physical parameters $\lambda_{i\ge1}$ to reproduce
\emph{both} the perpendicular (physical) and the parallel (phase) parts of $-iH\ket{\psi}$. Since each
physical parameter generically carries some phase, the fit becomes a compromise that distorts the
physical motion. The pure-phase knob $\lambda_0$ absorbs the parallel part, freeing the physical
parameters to track only the genuine motion.
\end{remark}

\section{Eliminating $\lambda_0$: emergence of $Q_\psi$}\label{sec:project}

Run the index in \eqref{eq:vel}--\eqref{eq:EOM} over $i=0,1,\dots,N_p$ and write
$\ket{\dot\psi}=i\dot\lambda_0\ket{\psi}+\ket{\chi}$, with the physical part
$\ket{\chi}=\sum_{i\ge1}\dot\lambda_i\ket{\partial_i\psi}$. Define the phase overlaps
\begin{equation}
  b_i\equiv\im\braket{\partial_i\psi}{\psi}\quad(i\ge1),
  \qquad\text{so}\qquad
  \braket{\partial_i\psi}{\psi}=ib_i,\;\;\braket{\psi}{\partial_i\psi}=-ib_i,
  \label{eq:bdef}
\end{equation}
the last equalities following from the normalization identity \eqref{eq:normid}. A short computation
of the two non-constant terms of \eqref{eq:Lexp2} gives
\begin{align}
  \braket{\dot\psi}{\dot\psi}
    &=\dot\lambda_0^2-2\dot\lambda_0\!\sum_i b_i\dot\lambda_i+\braket{\chi}{\chi},\\
  -2\,\im\mel{\dot\psi}{H}{\psi}
    &=2\dot\lambda_0\langle H\rangle-2\,\im\mel{\chi}{H}{\psi}.
\end{align}
Minimizing over the phase velocity, $\partial L/\partial\dot\lambda_0=0$, yields
\begin{equation}
  \dot\lambda_0=\sum_i b_i\dot\lambda_i-\langle H\rangle.
  \label{eq:lam0}
\end{equation}
Substituting \eqref{eq:lam0} back, the three $\dot\lambda_0$-dependent terms collapse to
$-\dot\lambda_0^2=-\big(\sum_i b_i\dot\lambda_i-\langle H\rangle\big)^2$. Expanding and collecting the
quadratic and linear pieces in the physical velocities reproduces the linear system $M\dlam=V$, but
now with \emph{phase-corrected} coefficients,
\begin{equation}
  \boxed{\;M_{ij}=\re\braket{\partial_i\psi}{\partial_j\psi}-b_ib_j,
  \qquad
  V_i=\im\mel{\partial_i\psi}{H}{\psi}-\langle H\rangle\,b_i.\;}
  \label{eq:MVproj}
\end{equation}

\section{Connection to the quantum geometric tensor}\label{sec:qgt}

The corrections in \eqref{eq:MVproj} are exactly the insertion of the projector
$Q_\psi\equiv I-\ket{\psi}\bra{\psi}$. Using \eqref{eq:bdef},
\begin{equation}
  \braket{\partial_i\psi}{\psi}\braket{\psi}{\partial_j\psi}=(ib_i)(-ib_j)=b_ib_j\in\mathbb{R},
\end{equation}
so that
\begin{align}
  M_{ij}&=\re\!\big(\braket{\partial_i\psi}{\partial_j\psi}-\braket{\partial_i\psi}{\psi}\braket{\psi}{\partial_j\psi}\big)
        =\re\mel{\partial_i\psi}{Q_\psi}{\partial_j\psi},\\
  V_i&=\im\mel{\partial_i\psi}{(H-\langle H\rangle)}{\psi}
      =\im\mel{\partial_i\psi}{Q_\psi H}{\psi},
  \label{eq:Qform}
\end{align}
where in the last line $Q_\psi H\ket{\psi}=(H-\langle H\rangle)\ket{\psi}$.

\begin{remark}[The projected $M$ is the Fubini--Study metric]
$M_{ij}=\re\mel{\partial_i\psi}{Q_\psi}{\partial_j\psi}$ is the real part of the quantum geometric
tensor, equal to one quarter of the quantum Fisher information. It is invariant under reparametrizing
the local phase; the bare $\re\braket{\partial_i\psi}{\partial_j\psi}$ is not. Mechanically,
\textbf{the projector is the global phase integrated out}: appending $\lambda_0$ and eliminating it
is identical to inserting $Q_\psi$. The vector $Q_\psi\ket{\partial_i\psi}$ is ``the part of the
$i$-th tangent direction that actually changes the physical state.''
\end{remark}

\begin{remark}[Conditioning]
If a pure-phase direction $i\ket{\psi}$ is reachable by the physical parameters, the bare $M$ keeps
that nearly-null direction in its column space and is ill-conditioned; the projected $M$ removes it
($Q_\psi\, i\ket{\psi}=0$), stabilizing the trajectory. In practice $M$ is still frequently
singular or ill-conditioned and is regularized (Section~\ref{sec:integrate}).
\end{remark}

\newpage
\part{Algorithm and Worked Examples}

\section{Evaluating $M$ and $V$}\label{sec:eval}

For a product ansatz of rotation blocks $u(\lambda)=e^{i\lambda\sigma/2}$ (with $\sigma$ a Pauli
string), each block obeys $\partial u/\partial\lambda=\tfrac{i}{2}u\sigma$, so the tangent vector is
\begin{equation}
  \ket{\partial_i\psi}=f_i\,\hat U_i(\boldsymbol{\lambda})\ket{\psi_0},\qquad f_i=\tfrac{i}{2},
  \label{eq:struct}
\end{equation}
where $\hat U_i$ is the ansatz with the $i$-th block's generator $\sigma_i$ inserted. All entries of
$M$ and $V$ then reduce to overlaps of the form $\braket{0}{\hat U_i^\dagger \hat U_j\,0}$ and
$\mel{0}{\hat U_i^\dagger\sigma_p U}{0}$, with $H=\sum_p h_p\sigma_p$ its Pauli decomposition.

\begin{itemize}[leftmargin=1.4em]
  \item \textbf{On hardware:} these overlaps are measured by a Hadamard test with one ancilla
  prepared in $(\ket{0}+\ket{1})/\sqrt2$. The cost is $\mathcal{O}(N_p^2)$ circuit evaluations per
  step for $M$.
  \item \textbf{Classical emulation (small systems):} compute $\ket{\psi}$ and $\ket{\partial_i\psi}$
  as state vectors directly and form the inner products. The Fubini--Study metric
  $M_{ij}=\re\mel{\partial_i\psi}{Q_\psi}{\partial_j\psi}$ is returned by standard library routines
  (e.g.\ \texttt{qml.metric\_tensor}); the projector correction $-\langle H\rangle b_i$ on $V$ must
  be added by hand, or one simply evaluates $V_i=\im\mel{\partial_i\psi}{(H-\langle H\rangle)}{\psi}$.
\end{itemize}

\section{Time integration and regularization}\label{sec:integrate}

The equation of motion $M\dlam=V$ is integrated forward in time. With a first-order (Euler) stepper
and step $\delta t$,
\begin{equation}
  \boldsymbol{\lambda}(t+\delta t)=\boldsymbol{\lambda}(t)+\delta t\,M^{-1}V .
  \label{eq:euler}
\end{equation}
Because $M$ is often singular or ill-conditioned, it is regularized before inversion. A simple and
common choice is a determinant-gated ridge,
\begin{equation}
  M\;\longrightarrow\;M+\varepsilon\, I \quad\text{whenever}\quad \det M<\varepsilon,
  \label{eq:reg}
\end{equation}
with e.g.\ $\varepsilon=10^{-7}$. The regularization parameter is a numerical knob (it trades
stability against bias), not a physical input.

\begin{remark}[Error control]
Euler is first order, so the time-discretization error accumulates linearly in $\delta t$; decreasing
$\delta t$ (or using a higher-order Runge--Kutta integrator) systematically improves fidelity.
\end{remark}

\section{The algorithm}\label{sec:algorithm}

\noindent\fbox{\parbox{0.95\linewidth}{
\textbf{Variational quantum simulation of a quench.}
\begin{enumerate}[leftmargin=1.6em,itemsep=2pt]
  \item \textbf{Initial state (VQE).} Obtain $\boldsymbol{\lambda}_{\mathrm{opt}}$ by minimizing
  $\mel{\psi(\boldsymbol{\lambda})}{H_{\mathrm{init}}}{\psi(\boldsymbol{\lambda})}$ for the
  pre-quench Hamiltonian. Set $\boldsymbol{\lambda}(0)=\boldsymbol{\lambda}_{\mathrm{opt}}$ and
  switch on the quenched Hamiltonian $H$.
  \item \textbf{Per step ($t\to t+\delta t$):}
  \begin{enumerate}[label=(\alph*),leftmargin=1.6em]
    \item Evaluate $M(\boldsymbol{\lambda})$ and $V(\boldsymbol{\lambda})$ via
          \eqref{eq:Qform}.
    \item Regularize $M$ by \eqref{eq:reg} and solve $\dlam=M^{-1}V$.
    \item Euler update \eqref{eq:euler}: $\boldsymbol{\lambda}\!\mathrel{+}=\delta t\,\dlam$.
    \item Measure observables $\langle O\rangle(t)$ in the updated state.
  \end{enumerate}
  \item \textbf{Validate} against an exact reference: fidelity
  $F(t)=|\braket{\Psi_{\mathrm{exact}}(t)}{\psi(\boldsymbol{\lambda}(t))}|^2$, conservation of any
  symmetry charge, and the residual $L_{\min}(t)$ of \eqref{eq:residual}.
\end{enumerate}}}

\medskip
Using the \emph{same} ansatz for VQE and VQS lets the evolution start from
$\boldsymbol{\lambda}_{\mathrm{opt}}$ at constant depth throughout, in contrast to Trotter circuits
whose depth grows with $t$.

\section{Minimal model I: naive freezing vs.\ projected tracking}\label{sec:ex1}

Consider a single qubit with a one-parameter ansatz that rotates the azimuthal angle at a
\emph{fixed} polar angle $\theta_\ast$,
\begin{equation}
  \ket{\psi(\phi)}=e^{-i\phi Z/2}\big(\cos\tfrac{\theta_\ast}{2}\ket{0}+\sin\tfrac{\theta_\ast}{2}\ket{1}\big),
  \qquad H=\tfrac{\omega}{2}X .
\end{equation}
The driving force $X$ rotates the Bloch vector about $\hat x$, changing the polar angle, which the
ansatz \emph{cannot} follow---an instructive imperfect-ansatz situation. A direct computation gives
\[
  M^{\text{bare}}=\tfrac14,\qquad
  b_\phi=\tfrac12\cos\theta_\ast,\qquad
  M^{\text{proj}}=\tfrac14\sin^2\theta_\ast,
\]
and, crucially,
\[
  \mel{\partial_\phi\psi}{H}{\psi}=-\tfrac{\omega}{4}\sin\theta_\ast\sin\phi\;\in\mathbb{R}
  \;\Longrightarrow\;
  V^{\text{bare}}=\im\mel{\partial_\phi\psi}{H}{\psi}=0 .
\]
\textbf{Bare (no phase handling):} $\dot\phi=V^{\text{bare}}/M^{\text{bare}}=0$ --- the state freezes
and fails to move at all. The physical force is hidden: the bare matrix element came out purely real
because it is dominated by the phase overlap, and the genuine signal lives in the imaginary part only
\emph{after} the phase is removed.

\textbf{Projected (with $\lambda_0$):} $V^{\text{proj}}=V^{\text{bare}}-\langle H\rangle b_\phi
=-\tfrac{\omega}{4}\sin\theta_\ast\cos\theta_\ast\cos\phi$, hence
\begin{equation}
  \dot\phi=\frac{V^{\text{proj}}}{M^{\text{proj}}}=-\omega\cot\theta_\ast\cos\phi .
  \label{eq:ex1result}
\end{equation}
This is precisely the orthogonal projection of the exact Bloch precession
$\dot{\mathbf r}=\omega\,\hat x\times\mathbf r$ onto the achievable azimuthal direction. The two
agree to machine precision in numerics (Table~\ref{tab:ex1}).

\begin{table}[h]
\centering
\renewcommand{\arraystretch}{1.25}
\begin{tabular}{ccccc}
\toprule
$\phi$ & $\dot\phi$ (bare) & $\dot\phi$ (projected) & exact best-fit & $b_\phi$ \\
\midrule
$0.0$ & $\approx 0$ & $-1.000$ & $-1.000$ & $0.354$ \\
$1.0$ & $\approx 0$ & $-0.540$ & $-0.540$ & $0.354$ \\
$2.0$ & $\approx 0$ & $+0.416$ & $+0.416$ & $0.354$ \\
\bottomrule
\end{tabular}
\caption{Single-qubit model with $\theta_\ast=\pi/4$, $\omega=1$. The bare update freezes
($\dot\phi=0$); the projected update matches the exact geometric best-fit \eqref{eq:ex1result}.}
\label{tab:ex1}
\end{table}

\begin{remark}[Moral]
The phase knob does not add reach; it tells the algorithm not to score phase as error. The very fact
that the physical parameter carries phase ($b_\phi\neq0$) is what fools the bare fit---so quotienting
the phase out is mandatory, not optional.
\end{remark}

\section{Minimal model II: the eigenstate limit}\label{sec:ex2}

The cleanest test is an eigenstate $H\ket{\psi}=E\ket{\psi}$. Exact evolution is
$e^{-iEt}\ket{\psi}$: the ray is stationary and only the global phase advances, which should sit
entirely in $\lambda_0$.

\textbf{Projected.} Since $(H-\langle H\rangle)\ket{\psi}=(H-E)\ket{\psi}=0$, equation
\eqref{eq:Qform} gives $V_i\equiv0$, so every physical velocity vanishes; the phase equation
\eqref{eq:lam0} gives $\dot\lambda_0=-E$. The physical state is frozen and the phase ticks at the
eigenvalue---exactly correct, independent of the ansatz.

\textbf{Bare.} Here $V_i^{\text{bare}}=\im\mel{\partial_i\psi}{H}{\psi}=E\,b_i$, generally nonzero:
the bare method tries to move physical parameters to reproduce the unobservable phase rotation. Whether
this corrupts observables depends on the ansatz:
\begin{itemize}[leftmargin=1.4em]
  \item If the ansatz happens to contain a free pure-phase direction, that knob absorbs the spurious
  velocity and the physical parameters stay frozen. A two-qubit example with $H=\tfrac{\omega}{2}(Z_0+Z_1)$,
  reference $\ket{00}$, and a $Z_0+Z_1$ rotation acting as the phase knob gives bare velocities
  $(\dot\theta,\dot\phi)=(0,1)$ and unit fidelity for all $t$: the physical direction $\theta$ is
  protected. This vindicates the intuition that a parameter able to play the role of $\lambda_0$
  shields the physical sector.
  \item If the ansatz contains no free phase direction (as in Model~I), the spurious phase-chasing
  contaminates the physical parameters and the trajectory drifts or freezes incorrectly.
\end{itemize}

\begin{remark}[What $\lambda_0$ guarantees]
A real-parameterized ansatz generically does \emph{not} span $i\ket{\psi}$ at every point of the
trajectory. Adding $\lambda_0$ (equivalently inserting $Q_\psi$) provides a dedicated phase outlet at
every step, making the physical parameters provably immune to phase contamination---no reliance on the
ansatz being accidentally favourable.
\end{remark}

\newpage
\appendix
\section{Summary of key results}

\begin{center}
\renewcommand{\arraystretch}{1.5}
\begin{tabular}{lll}
\toprule
\textbf{Quantity} & \textbf{Symbol} & \textbf{Expression}\\
\midrule
McLachlan functional & $L$ & $\big\|\,\ket{\dot\psi}+iH\ket{\psi}\,\big\|^2$\\
Equation of motion & $M\dlam=V$ & normal equations of $L$\\
Metric (bare) & $M_{ij}$ & $\re\braket{\partial_i\psi}{\partial_j\psi}$\\
Force (bare) & $V_i$ & $\im\mel{\partial_i\psi}{H}{\psi}$\\
Metric (projected) & $M_{ij}$ & $\re\mel{\partial_i\psi}{Q_\psi}{\partial_j\psi}$\\
Force (projected) & $V_i$ & $\im\mel{\partial_i\psi}{(H-\langle H\rangle)}{\psi}$\\
Phase velocity & $\dot\lambda_0$ & $\sum_i b_i\dot\lambda_i-\langle H\rangle$\\
Phase overlap & $b_i$ & $\im\braket{\partial_i\psi}{\psi}$\\
Residual & $L_{\min}$ & $\langle H^2\rangle-V^{\!\top}M^{-1}V$\\
Euler update & $\boldsymbol{\lambda}(t{+}\delta t)$ & $\boldsymbol{\lambda}(t)+\delta t\,M^{-1}V$\\
Regularization & $M$ & $M+\varepsilon I$ if $\det M<\varepsilon$\\
\bottomrule
\end{tabular}
\end{center}

\section*{Logical structure of the derivation}
\addcontentsline{toc}{section}{Logical structure of the derivation}
\begin{center}
\renewcommand{\arraystretch}{1.3}
\begin{tabular}{rl}
\toprule
Step & Content\\
\midrule
1 & Restrict dynamics to manifold $\ket{\psi(\boldsymbol{\lambda})}$; velocity $=\sum\dot\lambda_i\ket{\partial_i\psi}$ \\
2 & Define McLachlan distance $\|\ket{\dot\psi}+iH\ket{\psi}\|^2$ \\
3 & Expand $\Rightarrow$ quadratic form $\dlam^{\!\top}M\dlam-2V^{\!\top}\dlam+\langle H^2\rangle$ \\
4 & Minimize over $\dlam$ $\Rightarrow$ normal equations $M\dlam=V$ \\
5 & Add global phase $\lambda_0$; eliminate it $\Rightarrow$ insert $Q_\psi$ \\
6 & Integrate $M\dlam=V$ in time with regularization \\
\bottomrule
\end{tabular}
\end{center}

\section{Notation conventions}

\begin{center}
\renewcommand{\arraystretch}{1.3}
\begin{tabular}{ll}
\toprule
\textbf{Symbol} & \textbf{Meaning}\\
\midrule
$\ket{\psi(\boldsymbol{\lambda})}$ & variational (ansatz) state\\
$\boldsymbol{\lambda},\dlam$ & real parameters and their time derivatives (velocities)\\
$\lambda_0$ & global-phase parameter; $\ket{\partial_0\psi}=i\ket{\psi}$\\
$\ket{\partial_i\psi}$ & tangent vector $\partial\ket{\psi}/\partial\lambda_i$\\
$H$ & Hamiltonian (post-quench); $H=\sum_p h_p\sigma_p$ in Pauli form\\
$\langle O\rangle$ & expectation $\mel{\psi}{O}{\psi}$; in particular $\langle H\rangle,\langle H^2\rangle$\\
$M_{ij}$ & Fubini--Study metric (real part of quantum geometric tensor)\\
$V_i$ & projected energy-gradient vector\\
$Q_\psi$ & projector $I-\ket{\psi}\bra{\psi}$\\
$b_i$ & phase overlap $\im\braket{\partial_i\psi}{\psi}=\tfrac12\langle G_i\rangle$\\
$\delta t,\;\varepsilon$ & integration step; regularization parameter\\
$F(t)$ & fidelity to the exact reference state\\
$L_{\min}$ & minimal McLachlan distance squared (residual)\\
\bottomrule
\end{tabular}
\end{center}

\end{document}
